% from ega4-2-fr.tex
\oldpage[IV-2]{149}

\begin{proposition}[6.7.7]
\label{IV.2.s6.7.7-fr}
Soient $k$ un corps, $X$ un $k$-préschéma localement noethérien, $x$ un
point de $X$; désignons par $\boldsymbol{Q}(k')$ la relation «$k'$ est une extension de
$k$, et $\boldsymbol{P}(X_{(k')},x')$ est vraie pour tout point $x'$ de $X_{(k')}$ dont
$x$ est la projection dans $X$», où $\boldsymbol{P}(Z,z)$ est une des propriétés \textup{a),
b), c)} de l'énoncé \textup{(6.7.4)}. Alors les propriétés suivantes sont équivalentes:
\begin{enumerate}
\item[\textup{a)}] $\boldsymbol{Q}(k')$ est vraie pour toute extension finie $k'$ de $k$.
\item[\textup{b)}] $\boldsymbol{Q}(k')$ est vraie pour toute extension finie radicielle $k'$ de $k$.
\item[\textup{c)}] $\boldsymbol{Q}(k')$ est vraie pour toute extension de type fini $k'$ de $k$.
\end{enumerate}
Supposons en outre que $X$ soit localement de type fini sur $k$; alors les
trois propriétés précédentes sont aussi équivalentes aux suivantes:
\begin{enumerate}
\item[\textup{d)}] $\boldsymbol{Q}(k')$ est vraie pour toute extension $k'$ de $k$.
\item[\textup{e)}] $\boldsymbol{Q}(k')$ est vraie pour \textup{une} extension parfaite $k'$ de $k$.
\item[\textup{f)}] $\boldsymbol{Q}(k')$ est vraie pour toute extension $k'=k^{p^{-s}}$, où $p$ est
l'exposant caractéristique de $k$, et $s>0$.
\end{enumerate}
\end{proposition}

\begin{proof}
Pour prouver l'équivalence de \emph{a)}, \emph{b)} et \emph{c)}, il suffit évidemment d'établir
que \emph{b)} entraîne \emph{c)}. Soit donc $K$ une extension de type fini de $k$, qui est
par suite le corps des fractions d'une $k$-algèbre de type fini $A$; posons
$Y=\operatorname{Spec}(A)$. On sait (4.6.6) qu'il existe une extension finie
et radicielle $k'$ de $k$ telle que $Y'=(Y\otimes_k k')_{\mathrm{red}}$ soit
séparable sur $k'$; si $\eta'$ est le point générique de $Y'$, $K'=\boldsymbol{k}(\eta')$
est une extension séparable de $k'$ par (4.6.1). Cela étant,
$\boldsymbol{P}(X_{(k')},x')$ est vraie par hypothèse pour tout point $x'$ de $X_{(k')}$
au-dessus de $x$, donc il résulte de (6.7.4) que $\boldsymbol{P}(X_{(K')},x'')$ est vraie
pour tout point $x''$ de $X_{(K')}$ au-dessus de $x$, puisque
$X_{(K')}=(X_{(k')})_{(K')}$, que $x''$ est au-dessus d'un point $x'$ de
$X_{(k')}$, lui-même au-dessus de $x$, et que $K'$ est séparable sur $k'$.
Mais on a aussi
$X_{(K')}=(X_{(K)})_{(K')}$ et pour tout $x_0\in X_{(K)}$ au-dessus de
$x$, il existe $x''\in X_{(K')}$ au-dessus de $x_0$; il résulte donc de
(6.7.4) que $\boldsymbol{P}(X_{(K)},x_0)$ est vraie.

Supposons maintenant $X$ localement de type fini sur $k$, de sorte que pour
\emph{toute} extension $K$ de $k$, $X_{(K)}$ est localement noethérien. Comme une
extension radicielle de $k$ est isomorphe à une sous-extension d'une extension
parfaite quelconque de $k$, il résulte aussitôt de (6.7.4) que \emph{e)} implique \emph{f)};
de même, toute extension finie radicielle de $k$ est contenue dans une
extension $k^{p^{-s}}\subset k^{p^{-\infty}}$, donc \emph{f)} entraîne \emph{b)}; \emph{d)} entraîne
trivialement \emph{e)}, et enfin \emph{e)} entraîne \emph{d)}. En effet, si $K$ est une extension
\emph{quelconque} de $k$,
il existe une extension $K'$ de $k$ contenant (à des $k$-isomorphismes près)
$k'$ et $K$; comme $K'$ est extension \emph{séparable} de $k'$, on déduit de (6.7.4)
que $\boldsymbol{Q}(k')$ est vraie, puis que $\boldsymbol{Q}(K)$ est vraie, en raisonnant comme dans la
première partie de la démonstration. Il suffit donc de prouver que \emph{b) entraîne
e)}. Nous prendrons $k'=k^{p^{-\infty}}$ dans \emph{e)}.

La question étant locale sur $X$, on peut en outre se borner au cas où $X$
est affine et de type fini sur $k$, en remplaçant $X$ par un voisinage de $x$;
soit donc $X=\operatorname{Spec}(B)$, où $B$ est une $k$-algèbre de type fini;
en outre, par définition de la propriété $\boldsymbol{P}$ dans les cas considérés, on peut
aussi remplacer $X$ par le schéma local de $X$ au point $x$, donc supposer que
$B$ est un anneau \emph{local} noethérien. Posons $B'=B\otimes_k k'$; $k'$ est
limite inductive de ses sous-extensions \emph{finies} $k'_\lambda$ et si l'on pose
$B'_\lambda=B\otimes_k k'_\lambda$, le morphisme
$f_\lambda:\operatorname{Spec}(B')\to\operatorname{Spec}(B'_\lambda)$ est
un \emph{homéomorphisme} de $\operatorname{Spec}(B')$ \emph{sur}
$\operatorname{Spec}(B'_\lambda)$ pour tout $\lambda$.

